\documentclass[11pt]{article}

\usepackage{amsmath,amssymb}
\usepackage{xurl}
\usepackage[hidelinks]{hyperref}
\setlength{\emergencystretch}{2em}

% Shared commands for notes in docs/paper-gaps/.

\DeclareUnicodeCharacter{2080}{\ifmmode{}_{0}\else\textsubscript{0}\fi}
\DeclareUnicodeCharacter{2081}{\ifmmode{}_{1}\else\textsubscript{1}\fi}
\DeclareUnicodeCharacter{2082}{\ifmmode{}_{2}\else\textsubscript{2}\fi}
\DeclareUnicodeCharacter{2099}{\ifmmode{}_{n}\else\textsubscript{n}\fi}

\newcommand{\Lean}{\textsc{Lean}}
\ExplSyntaxOn
\NewDocumentCommand{\leanid}{m}{
  \ifmmode
    \textsf{\detokenize{#1}}
  \else
    \group_begin:
    \urlstyle{sf}
    \tl_set:Nn \l_tmpa_tl {#1}
    \tl_replace_all:Nnn \l_tmpa_tl {\_} {_}
    \exp_args:NV \path \l_tmpa_tl
    \group_end:
  \fi
}
\ExplSyntaxOff
\providecommand{\tr}{\operatorname{tr}}
\newcommand{\tnleanIssueBase}{https://github.com/LionSR/TNLean/issues/}
\newcommand{\tnleanPrBase}{https://github.com/LionSR/TNLean/pull/}
\newcommand{\tnleanDocBase}{https://sirui-lu.com/TNLean/docs/}
\newcommand{\ghissue}[1]{\href{\tnleanIssueBase#1}{GitHub issue \##1}}
\newcommand{\ghpr}[1]{\href{\tnleanPrBase#1}{PR \##1}}
\newcommand{\doclink}[1]{\href{\tnleanDocBase#1.html}{\path{#1}}}

% Dirac notation and the identity operator, as in blueprint/src/macros/common.tex.
% \providecommand keeps note-local definitions authoritative.
\usepackage{dsfont}
\providecommand{\ket}[1]{|#1\rangle}
\providecommand{\bra}[1]{\langle #1|}
\providecommand{\braket}[2]{\langle #1 | #2 \rangle}
\providecommand{\ketbra}[2]{|#1\rangle\!\langle #2|}
\providecommand{\Id}{\mathds{1}}
\providecommand{\one}{\mathds{1}}
\providecommand{\C}{\mathbb{C}}
\providecommand{\MN}[1]{M_{#1}(\C)}
\providecommand{\MD}{M_D(\C)}

% External referencing for paper sources.  The build helper writes sanitized
% auxiliary files under build/paper-aux/ and excludes duplicated source labels.
\usepackage{xr}

% Import a generated paper auxiliary file with a caller-supplied label prefix.
% The candidate paths support builds from docs/paper-gaps/, the repository root,
% and build/paper-aux/ itself.
\newcommand{\importpaperaux}[2]{%
  \IfFileExists{../../build/paper-aux/#2.aux}{%
    \externaldocument[#1]{../../build/paper-aux/#2}%
  }{%
    \IfFileExists{build/paper-aux/#2.aux}{%
      \externaldocument[#1]{build/paper-aux/#2}%
    }{%
      \IfFileExists{#2.aux}{%
        \externaldocument[#1]{#2}%
      }{}%
    }%
  }%
}

% General external reference macros
\newcommand{\paperref}[2]{\ref{#1-#2}}
\newcommand{\papereqref}[2]{(\ref{#1-#2})}

% CPSV16 source (1606.00608 / MPDO-22-12-17-2.tex) external document and shortcuts
\importpaperaux{cpsv16-}{1606.00608/MPDO-22-12-17-2-tnlean}
\newcommand{\cpsvref}[1]{\paperref{cpsv16}{#1}}
\newcommand{\cpsveqref}[1]{\papereqref{cpsv16}{#1}}

% Verdict marker: \gapnote{<kind>}{<status>}. Kind is the policy
% classification (clarification, local-correction, scope-restriction,
% unfaithful, false-source, open-gap); status is open, wip (elimination
% underway), resolved, or historical. Machine-read by the site index;
% prints a verdict line.
\newcommand{\gapnote}[2]{\par\noindent\textbf{Verdict:} #1 (#2).\par}


\title{Common Sector Relabeling as an Explicit Hypothesis}
\author{}
\date{2026-05-10}

\begin{document}
\maketitle

\section{The assertion}

Several after-blocking sector-comparison theorems in the formalization take an explicit
hypothesis \leanid{CommonSectorRelabelingHypothesis} of type \texttt{Prop}.  The hypothesis
asserts: for every common blocked cyclic-sector family \(F\) with blocking period \(p\) and
block families \((A_k)_k\), the weighted block-diagonal tensor assembled from the
physically blocked blocks agrees, as an MPV family, with the same family assembled after
the explicit blocked-word relabeling carried by \(F\):
\[
  \forall F,\quad
  \bigoplus_k \mu_k^p \cdot A_k^{[p]} \sim_{\mathrm{MPV}}
  \bigoplus_k \mu_k^p \cdot R_F(A_k^{[p]}).
\]
Here \(A_k^{[p]}\) denotes the length-\(p\) blocking of the \(k\)-th primitive block, and
\(R_F(A_k^{[p]})\) is the same tensor read through the explicit word relabeling of
blocked physical indices.\footnote{%
  In the formal statement this relabeled block is the field
  \leanid{CommonBlockedCyclicSectorFamily.commonReindexedBlock}.}

\section{Relation to the paper}

The paper arXiv:1606.00608 does not state this as a separate hypothesis.
In the source argument the blocked-word identification is an automatic consequence of the
canonical blocking construction: blocking a tensor and then reading the result through the
explicit multi-index grouping of \(\{1,\ldots,d\}^p\) is definitionally the same operation.
The paper therefore treats the equality as transparent and does not name it.

\section{Status in the formalization}

The hypothesis is always provable.  The Lean declaration
\leanid{CommonGroupedBlockCastHypothesis.of\_flattenWordOfBlock\_cast\_eq} proves the
stronger coordinate-grouping assertion \leanid{CommonGroupedBlockCastHypothesis}
unconditionally, and
\leanid{CommonGroupedBlockCastHypothesis.toRelabelingHypothesis}
derives \leanid{CommonSectorRelabelingHypothesis} from it.%
\footnote{%
  Both declarations are in
  \doclink{TNLean/MPS/CanonicalForm/SectorComparison/CommonSectorTransport.lean}.
  The after-blocking sector-comparison theorems are in
  \doclink{TNLean/MPS/CanonicalForm/SectorComparison/ProportionalComparison.lean}.
}

The explicit parameterization is therefore a \emph{formalization design choice}, not a
genuine extra mathematical assumption.  Taking the hypothesis as an argument
exposes the logical boundary between the common-sector structural theorem and the
downstream comparison, making the dependency explicit at the type level.  Every caller can
discharge it with the one-line proof
\texttt{CommonGroupedBlockCastHypothesis.of\_flattenWordOfBlock\_cast\_eq d
  |>.toRelabelingHypothesis}.

\section{Verdict}

No deviation from the source paper: the hypothesis is always true and the proof is
available.  The formalization makes an implicit paper convention explicit at the
type level; this is a recognized formalization pattern, not a scope restriction or
an error.  No corrective action is required, but callers of the parameterized theorems
should discharge this hypothesis using the available lemma rather than treating it as
an open obligation.

\bibliographystyle{alpha}
\bibliography{references}

\end{document}
